\begin{figure}[tbp]
\centering
\begin{tikzpicture}[x=1.45cm,y=1cm,>=Stealth,
  every node/.style={font=\small}]
  \foreach \k in {0,...,5} {
    \node (v\k) at (\k,1.1) {$V_{\k}$};
  }
  \foreach \k in {0,...,4} {
    \pgfmathtruncatemacro{\next}{\k+1}
    \draw[->] (v\next.west) -- (v\k.east);
  }
  \draw[->] (v0.north) to[bend left=35]
    node[above] {$\Phi$} (v2.north);
  \draw[->] (v0.north) to[bend left=40]
    node[above] {$[1]$} (v5.north);
  \foreach \k in {0,...,4} {
    \fill[blue!6] (\k,-0.05) rectangle (\k+1,0.4);
    \node at (\k+0.5,0.175) {$\mathcal S_{\k}$};
  }
  \draw (0,-0.05) -- (5,-0.05);
  \foreach \k in {0,...,5} {
    \draw (\k,0.4) -- (\k,-0.12);
  }
  \node[below] at (0,-0.12) {$0$};
  \foreach \k in {1,...,4} {
    \node[below] at (\k,-0.12) {$\k/5$};
  }
  \node[below] at (5,-0.12) {$1$};
  \node[below] at (2.5,-0.65) {charge phase};
\end{tikzpicture}
\caption{The fifth-phase grid. Inclusion arrows point toward the larger
aisle. The autoequivalence $\Phi$ advances two indices, and the shift $[1]$
advances five. A nonzero object of $\mathcal S_k$ has charge phase in the
closed interval $[k/5,(k+1)/5]$, by Proposition~\ref{prop:block-support}.
At a shared endpoint, generators from both adjacent blocks enter the
definition of the slicing. The diagram shows the unperturbed grid;
Proposition~\ref{prop:fine-grid} gives the independent-parameter nesting.}
\label{fig:fifth-phase-grid}
\end{figure}
